\documentclass[11pt,a4paper,uplatex,dvipdfmx]{jsarticle}

\usepackage[margin=22mm]{geometry}
\usepackage{longtable}
\usepackage{booktabs}
\usepackage{array}
\usepackage{xcolor}
\usepackage{hyperref}
\usepackage{listings}
\usepackage{enumitem}

\hypersetup{
  colorlinks=true,
  linkcolor=blue,
  urlcolor=blue
}

\lstdefinestyle{aipl}{
  basicstyle=\ttfamily\small,
  breaklines=true,
  frame=single,
  columns=fullflexible,
  keepspaces=true
}

\setlist[itemize]{leftmargin=2em}
\setlist[enumerate]{leftmargin=2em}
\renewcommand{\arraystretch}{1.2}

\title{Python版AIPL ユーザーズガイド}
\author{児玉靖司（Yasushi Kodama）\\法政大学（Hosei University）}
\date{2026年5月12日}

\begin{document}
\maketitle
\tableofcontents
\newpage

\section{概要}

Python版AIPLは、\texttt{src/python-aipl/} にある研究・拡張用ランタイムである。
OCaml版よりも新しい言語機能を多く持ち、型注釈、効果注釈、所有権、線形型、
構造化並行、動的コンパイル、AI呼び出し、モデル検査などを扱える。

Python版AIPLは単なるABCL風アクター処理系ではなく、型付きアクター言語、
AI統合言語、分散アクター実験基盤、モデル検査対象言語、動的自己拡張言語
として利用できる。

\section{対象実装}

\begin{longtable}{p{0.30\linewidth}p{0.62\linewidth}}
\toprule
項目 & ファイル \\
\midrule
ランタイム本体 & \texttt{src/python-aipl/aipl\_interp.py} \\
パーサ/構文 & \texttt{src/python-aipl/grammar.lark}, \texttt{src/python-aipl/aipl\_parser.py} \\
AST & \texttt{src/python-aipl/aipl\_ast.py} \\
型検査 & \texttt{src/python-aipl/aipl\_typeck.py} \\
実行入口 & \texttt{src/python-aipl/aipl\_main.py} \\
サンプル & \texttt{src/python-aipl/samples/}, \texttt{samples-ai/}, \texttt{samples-remote/}, \texttt{samples-mc/} \\
\bottomrule
\end{longtable}

\section{全体構成}

\begin{center}
\begin{tabular}{c}
\fbox{\texttt{.abcl source}} \\
$\downarrow$ \\
\fbox{\texttt{grammar.lark / aipl\_parser.py}} \\
$\downarrow$ \\
\fbox{\texttt{aipl\_ast.py}} \\
$\downarrow$ \\
\fbox{\texttt{aipl\_typeck.py}} \\
$\downarrow$ \\
\fbox{\texttt{aipl\_interp.py actor runtime}} \\
$\downarrow$ \\
\fbox{AI / Remote / File / Image / Model Check integrations}
\end{tabular}
\end{center}

\section{基本言語機能}

Python版AIPLは、\texttt{class}、\texttt{method}、\texttt{var}、\texttt{function}
を中心にしたアクター指向言語である。クラスから生成されたインスタンスは
アクターとして動作し、\texttt{send}、\texttt{now}、\texttt{future}、
\texttt{await} によって通信する。

\begin{longtable}{p{0.34\linewidth}p{0.56\linewidth}}
\toprule
構文 & 意味 \\
\midrule
\texttt{send actor.method(args);} & 非同期送信。呼び出し元は待たない。 \\
\texttt{now actor.method(args)} & 同期呼び出し。相手の \texttt{reply} を待つ。 \\
\texttt{future actor.method(args)} & Futureを返す非同期呼び出し。 \\
\texttt{await(f)} / \texttt{await f} & Futureの結果を待つ。 \\
\texttt{self} & 現在のアクター自身。 \\
\texttt{sender} & メッセージ送信元。 \\
\texttt{become Class(args);} & 現在のアクターの振る舞いを別クラスへ変更する。 \\
\bottomrule
\end{longtable}

\section{型システム}

Python版AIPLは、OCaml版より新しい型付き構文を持つ。変数、引数、戻り値に
型注釈を書ける。

\begin{lstlisting}[style=aipl]
method add(a: int, b: int) -> int {
  return a + b;
}
\end{lstlisting}

対応する主な型を以下に示す。

\begin{itemize}
  \item \texttt{int}
  \item \texttt{float}
  \item \texttt{string}
  \item \texttt{bool}
  \item \texttt{array[T]}
  \item \texttt{tuple(...)}
  \item \texttt{record\{...\}}
  \item Union型
  \item 型変数
  \item 長さ付き配列
  \item \texttt{linear T}
\end{itemize}

\texttt{typeof(x)} により実行時の構造型を確認できる。

\section{Phase 11から17の研究機能}

\begin{longtable}{p{0.20\linewidth}p{0.70\linewidth}}
\toprule
Phase & 機能 \\
\midrule
Phase 11 & 型注釈、Union型、Generics、長さ付き配列。 \\
Phase 12 & capability effects。例: \texttt{!\{fs, ai, net, mut\}}。 \\
Phase 13 & CSP風チャンネル。 \\
Phase 14 & \texttt{linear} による線形値、use-after-move検出。 \\
Phase 15 & \texttt{pub} フィールドと所有権、外部書き込み禁止。 \\
Phase 16 & transient cast、\texttt{any -> T} 境界の実行時検査。 \\
Phase 17 & \texttt{scope \{ ... \}} による構造化並行。 \\
\bottomrule
\end{longtable}

\section{データ構造}

配列、タプル、レコードを扱える。配列は \texttt{[1, 2, 3]} のような
リテラル、\texttt{array\_push}、\texttt{array\_get}、\texttt{array\_set}
などで操作する。多次元配列もあり、\texttt{var grid[rows][cols];} のように
動的サイズで宣言できる。

タプルは \texttt{(1, "a")}、レコードは
\texttt{\{ name: "Alice", age: 30 \}} のように書ける。レコードは
\texttt{r.name} のようなドットアクセスに対応している。

\section{関数}

トップレベル関数とクラス内関数を定義できる。関数は \texttt{return} によって
同期的に値を返す。クラス内関数は、そのクラスのメソッドから補助関数として
呼び出せる。

\texttt{typeof(function\_name)} により、実行時に観測された呼び出し履歴から
推定された関数シグネチャを表示できる。

\section{アクター機能}

アクターはメッセージキューを持ち、非同期に動作する。
\texttt{reply(value)} によって \texttt{now} や \texttt{future} の呼び出し元に
値を返す。

\texttt{become Class(args);} によって、実行中のアクターの振る舞いを別クラスへ
変更できる。

\section{チャンネル}

CSP風のチャンネル機能がある。アクター間通信とは別に、明示的な同期キューを
使える。

\begin{lstlisting}[style=aipl]
channel(capacity, type_name)
channel_send(ch, value)
channel_recv(ch)
channel_try_recv(ch)
channel_close(ch)
channel_size(ch)
select_recv([ch1, ch2], timeout_ms)
\end{lstlisting}

\section{動的コンパイルと動的生成}

\texttt{compile(source)} により、文字列として渡したAIPLソースを実行中に
コンパイルできる。\texttt{spawn(class\_name, args...)} により、文字列で
指定したクラスからアクターを生成できる。

さらに、以下のメソッドインジェクション機能がある。

\begin{lstlisting}[style=aipl]
add_method(target, source)
remove_method(target, name)
methods_of(target)
\end{lstlisting}

これにより、クラスまたは特定アクターへメソッドを動的に追加・削除できる。

\section{AI連携}

Python版AIPLは、生成AI呼び出しを組み込み関数として持っている。

\begin{lstlisting}[style=aipl]
ai_call(prompt)
ai_call_with_system(system, prompt)
ai_call_priority(priority, prompt)
ai_call_retry(max_attempts, prompt)
ai_call_image(prompt, image)
ai_usage()
ai_remaining()
ai_cost()
\end{lstlisting}

Gemini、Anthropic Claude、OpenAIを切り替えられる。
\texttt{AIPL\_AI\_PROVIDER=mock} を使うと、外部APIなしでテストできる。

また、起動時に \texttt{AI} アクターが自動生成され、
\texttt{now AI.ask("...")} や \texttt{future AI.ask("...")} の形でAIを
アクターとして扱える。

\section{ファイル、JSON、画像処理}

Python版AIPLには、アプリ生成やデータ処理向けの組み込みもある。

\begin{longtable}{p{0.30\linewidth}p{0.60\linewidth}}
\toprule
関数 & 概要 \\
\midrule
\texttt{read\_file} & テキストファイルを読む。 \\
\texttt{write\_file} & テキストファイルを書く。 \\
\texttt{append\_file} & テキストファイルへ追記する。 \\
\texttt{read\_bytes} & バイナリを読む。 \\
\texttt{write\_bytes} & バイナリを書く。 \\
\texttt{json\_parse} & JSON文字列をAIPL値へ変換する。 \\
\texttt{json\_stringify} & AIPL値をJSON文字列へ変換する。 \\
\texttt{image\_load} & 画像を読み込む。 \\
\texttt{image\_save} & 画像を保存する。 \\
\texttt{image\_create} & 画像を生成する。 \\
\texttt{image\_pixel} & ピクセル値を読む。 \\
\texttt{image\_set\_pixel} & ピクセル値を書き込む。 \\
\texttt{image\_size} & 画像サイズを返す。 \\
\bottomrule
\end{longtable}

画像処理はPillowを使う。AIPLだけでHTML、CSS、画像、JSON manifestを生成する
サンプルもある。

\section{分散アクター}

HTTP経由で別プロセス、別ランタイムのアクターと通信できる。

\begin{lstlisting}[style=aipl]
web_listen(port)
web_expose(name, actor)
remote_call(host, actor, method, args...)
remote_now(host, actor, method, args...)
remote_future(host, actor, method, args...)
serve_forever()
\end{lstlisting}

OCaml版との相互通信も想定されており、JSON形式の
\texttt{/api/json/send} と \texttt{/api/json/call} を使う。

\section{モデル検査}

Python版AIPLには、アクタープログラムの bounded model checking 用モジュールがある。
\texttt{aipl\_modelcheck.py}、\texttt{aipl\_modelcheck\_load.py} により、
食事する哲学者やロック順序などのサンプルを検査できる。

また、Promela/SPIN向け出力とTLA+向け出力もある。

\begin{itemize}
  \item \texttt{aipl\_to\_promela.py}
  \item \texttt{aipl\_to\_tla.py}
  \item \texttt{samples-mc/*.abcl}
\end{itemize}

\section{実行方法}

Python版AIPLのメイン実行ファイルは次である。

\begin{lstlisting}[style=aipl]
src/python-aipl/aipl_main.py
\end{lstlisting}

この環境ではラッパーとして次も使っている。

\begin{lstlisting}[style=aipl]
scripts/run_aipl_python313.sh
\end{lstlisting}

実行例:

\begin{lstlisting}[style=aipl,language=bash]
cd /Users/yaskodama/local-genai-chatgpt-ga/aios-claude
scripts/run_aipl_python313.sh src/python-aipl/samples/Hello.abcl
\end{lstlisting}

\section{サンプルとテスト}

サンプルは主に以下にある。

\begin{itemize}
  \item \texttt{src/python-aipl/samples/}
  \item \texttt{src/python-aipl/samples-ai/}
  \item \texttt{src/python-aipl/samples-remote/}
  \item \texttt{src/python-aipl/samples-mc/}
\end{itemize}

テストは \texttt{\_test\_*.py} として多数用意されており、型検査、配列、
タプル、レコード、チャンネル、線形型、所有権、メソッドインジェクション、
AIアクター、モデル検査などを確認できる。

\section{まとめ}

Python版AIPLは、型付きアクター言語としての中核に加え、AI統合、動的拡張、
分散アクター、モデル検査、ファイル・画像・JSON処理を統合した実験用言語である。
OCaml版がネイティブ実行とアクターランタイムの基礎を担う一方で、Python版は
新しい言語仕様とAI連携を素早く検証するための中心的な実装である。

\end{document}
